\documentclass[tikz,border=10pt]{standalone}
\usepackage{tikz}
\usetikzlibrary{shapes.geometric, arrows.meta, positioning, fit, backgrounds, calc}
\usepackage{xeCJK}
\setCJKmainfont{SimSun}

\begin{document}
\begin{tikzpicture}[
    node distance=1.2cm and 1.5cm,
    >={Stealth[length=3mm]},
    box/.style={rectangle, draw, rounded corners, minimum height=1cm, minimum width=2.8cm, align=center, font=\small},
    inputbox/.style={box, fill=blue!15, draw=blue!60},
    methodbox/.style={box, fill=green!15, draw=green!60},
    spotbox/.style={box, fill=orange!15, draw=orange!60},
    ensemblebox/.style={box, fill=purple!15, draw=purple!60},
    outputbox/.style={box, fill=red!15, draw=red!60},
    modulelabel/.style={font=\small\bfseries, align=center},
    arrow/.style={->, thick, draw=gray!70}
]

% 输入层
\node[inputbox, minimum width=4cm] (input) {多节点监控时序数据\\$\{x_1^t, x_2^t, \ldots, x_N^t\}$};

% 多视角相似性度量模块
\node[methodbox, below left=1.5cm and 0.5cm of input] (euclidean) {欧氏距离\\数值相似性};
\node[methodbox, below=1.5cm of input] (autoencoder) {自编码器嵌入\\形态相似性};
\node[methodbox, below right=1.5cm and 0.5cm of input] (dbscan) {DBSCAN聚类\\密度相似性};

% 异常分数
\node[box, fill=gray!10, below=0.8cm of euclidean, minimum width=2cm] (score1) {异常分数 $s_1^t$};
\node[box, fill=gray!10, below=0.8cm of autoencoder, minimum width=2cm] (score2) {异常分数 $s_2^t$};
\node[box, fill=gray!10, below=0.8cm of dbscan, minimum width=2cm] (score3) {异常分数 $s_3^t$};

% I-SPOT自适应阈值模块
\node[spotbox, below=0.8cm of score1, minimum width=2cm] (spot1) {I-SPOT\\自适应阈值};
\node[spotbox, below=0.8cm of score2, minimum width=2cm] (spot2) {I-SPOT\\自适应阈值};
\node[spotbox, below=0.8cm of score3, minimum width=2cm] (spot3) {I-SPOT\\自适应阈值};

% 判定结果
\node[box, fill=gray!10, below=0.8cm of spot1, minimum width=1.8cm] (decision1) {$s_1^t > z_1^t$};
\node[box, fill=gray!10, below=0.8cm of spot2, minimum width=1.8cm] (decision2) {$s_2^t > z_2^t$};
\node[box, fill=gray!10, below=0.8cm of spot3, minimum width=1.8cm] (decision3) {$s_3^t > z_3^t$};

% 集成决策模块
\node[ensemblebox, below=1.2cm of decision2, minimum width=5cm] (ensemble) {保守集成决策（逻辑与）\\$\hat{y}_t = \mathbb{I}[\bigwedge_{m=1}^{M}(s_t^{(m)} > z_t^{(m)})]$};

% 输出
\node[outputbox, below=1cm of ensemble, minimum width=3.5cm] (output) {异常检测结果\\节点$i$时刻$t$异常判定};

% 连接箭头
\draw[arrow] (input) -- (euclidean);
\draw[arrow] (input) -- (autoencoder);
\draw[arrow] (input) -- (dbscan);

\draw[arrow] (euclidean) -- (score1);
\draw[arrow] (autoencoder) -- (score2);
\draw[arrow] (dbscan) -- (score3);

\draw[arrow] (score1) -- (spot1);
\draw[arrow] (score2) -- (spot2);
\draw[arrow] (score3) -- (spot3);

\draw[arrow] (spot1) -- (decision1);
\draw[arrow] (spot2) -- (decision2);
\draw[arrow] (spot3) -- (decision3);

\draw[arrow] (decision1) -- (ensemble);
\draw[arrow] (decision2) -- (ensemble);
\draw[arrow] (decision3) -- (ensemble);

\draw[arrow] (ensemble) -- (output);

% 模块标签背景框
\begin{scope}[on background layer]
    % 多视角相似性度量模块
    \node[draw=green!60, dashed, rounded corners, fit=(euclidean)(autoencoder)(dbscan), inner sep=8pt, fill=green!5] (module1) {};
    \node[modulelabel, above=0.1cm of module1.north, text=green!70!black] {多视角相似性度量模块};

    % 自适应阈值调整模块
    \node[draw=orange!60, dashed, rounded corners, fit=(spot1)(spot2)(spot3), inner sep=8pt, fill=orange!5] (module2) {};
    \node[modulelabel, right=0.3cm of module2.east, text=orange!70!black, rotate=0] {自适应阈值\\调整模块};

    % 集成决策模块
    \node[draw=purple!60, dashed, rounded corners, fit=(ensemble), inner sep=8pt, fill=purple!5] (module3) {};
    \node[modulelabel, right=0.3cm of module3.east, text=purple!70!black] {集成决策模块};
\end{scope}

\end{tikzpicture}
\end{document}
